\documentclass[../../../include/open-logic-section]{subfiles}

\begin{document}

\olfileid{sth}{story}{urelements}
\olsection{غیر سیٹ بنیادی شےواں منّیئے؟}

اگلے کجھ باباں وچ اسیں سیٹ دے اکٹھا کرن والے مرحلہ وار تصور توں مسلّمے
کڈھن دی کوشش کراں گے۔ پر اگے ودھن توں پہلاں \emph{غیر سیٹ بنیادی شےواں}
بارے کجھ ہور کہنا پئے گا۔

\olref[approach]{sec} دی صورت سانوں نویں سیٹ صرف پرانے \emph{سیٹاں}
توں بناؤن دی اجازت دیندی سی۔ پر ہو سکدا اے اسیں منّنا چاہیئے کہ کجھ
\emph{غیر سیٹ}، گایاں، سُور، ریت دے ذرے یا ہور کجھ، سیٹاں دے !!{element}s
ہو سکن۔ اوس صورت وچ اسیں کجھ بنیادی شےواں، \emph{غیر سیٹ بنیادی شےواں}،
نال شروع کراں گے۔ فیر آکھاں گے کہ ہر مرحلے $S$ اُتے غیر سیٹ بنیادی شےواں
یا $S$ توں پہلاں بنے سیٹاں نوں (کسے وی میل نال) اکٹھا کر کے ’’سارے ممکن‘‘
سیٹ بناواں گے۔ نتیجے دی صورت کجھ ایہو جیہی ہووے گی:
\begin{center}
	\begin{tikzpicture}[scale=0.6]
	\tikzset{cut_here/.style={densely dotted}}
	\draw (-4,6) -- (-1, 0) -- (1,0) -- (4,6); 
	\draw[cut_here] (-5,8)--(-4,6);
	\draw[cut_here] (5,8)--(4,6);
	\draw(-1.5, 1)--(1.5, 1);
	\draw(-2, 2)--(2, 2);
	\draw(-2.5, 3)--(2.5,3);
	\draw(-3, 4)--(3, 4);
	\draw(-3.5, 5)--(3.5, 5);
	\draw(-4, 6)--(4, 6);
	\node[label] at (2, 0) {\small 0};
	\node[label] at (2.5, 1) {\small 1};
	\node[label] at (3, 2) {\small 2};
	\node[label] at (3.5, 3) {\small 3};
	\node[label] at (4, 4) {\small 4};
	\node[label] at (4.5, 5) {\small 5};
	\node[label] at (5, 6) {\small 6};
	\end{tikzpicture}
\end{center}
سو ہُن سانوں فیصلہ کرنا اے۔ \emph{کی اسیں غیر سیٹ بنیادی شےواں منّیئے؟}

فلسفیانہ لحاظ نال اپنے نظریے وچ غیر سیٹ بنیادی شےواں شامل کرن دی تُک بندی
اے۔ وڈی وجہ ایہ اے کہ سیٹ نظریہ \emph{لاگو کرن جوگا} ہووے۔ مثال لئی
\olref[sfr][siz][]{chap} توں یاد کرو کہ دو سیٹ $A$ تے~$B$ اکّو جِڈّے،
یعنی $\cardeq{A}{B}$، اُدوں تے صرف اُدوں نیں جدوں اوہناں وچ اک یک بہ یک
تے پورا نقشہ ہووے۔ جے کھیت دیاں گایاں تے واڑے دے سُور، دوناں دے سیٹ ہون،
تاں اسیں ’’گایاں تے سُور اکّو گنتی وچ نیں‘‘ والے دعوے نوں سیٹ نظریے نال
بیان کر سکدے آں۔ پر جے اسیں غیر سیٹ بنیادی شےواں اُتے روک لائیے، تاں گایاں
تے سُوراں دے سیٹ \emph{نہیں} بنن گے تے اوہ بیان ممکن نہیں ہووے گا۔
اصل وچ سیٹ نظریہ سیٹاں توں چھٹ ہور کسے شے اُتے سیدھے طریقے نال لاگو ای
نہیں کر سکاں گے۔ (غیر سیٹ بنیادی شےواں شامل کرن دیاں ہور وجوہاں لئی
\citealt[pp.~vi, 24, 50--1]{Potter2004} ویکھو۔)

پر ریاضی وچ غیر سیٹ بنیادی شےواں منّن دا رواج بہت گھٹ اے۔ اک وجہ ایہ
اے کہ اوہناں توں بناں سیٹ نظریہ لکھنا \emph{ماڑا جیہا} سوکھا اے۔
پر کدی کدی اوہناں نوں سیٹ نظریے توں باہر رکھن لئی ہور دلچسپ دلیلاں ملدیاں نیں:
\begin{quote}
	جے اسیں منّ دے آں کہ سیٹ نظریہ ریاضی دی بنیاد اے، تاں سانوں صرف سیٹاں
	دی گل کر کے ساری ریاضی بیان کرن جوگا ہونا چاہیدا اے۔ سو ساڈے متغیر دی
	حد وچ گایاں تے سُوراں ورگیاں شےواں نہیں ہونیاں چاہیدیاں۔
	%But if $C$ is a cow, $\{C\}$ is a set, but not a legitimate mathematical object.
	\citep[p.~8]{Kunen1980}
\end{quote}
سو جے لاگو ہون اُتے زور دیئیے، تاں غیر سیٹ بنیادی شےواں \emph{شامل کرن}
دا خیال آؤندا اے۔ جے گھٹا کے بنیاد دین والا مقصد (ریاضی نوں خالص سیٹ
نظریے تک گھٹانا) سامھنے ہووے، تاں اوہناں نوں \emph{باہر رکھن} دا خیال آ
سکدا اے۔ ماڑی جیہی سُستی وی اوہناں نوں باہر رکھن ول اشارہ کردی اے۔

اسیں سبھ توں سُستا راہ پھڑاں گے۔ پر کچھ وجہ پڑھاؤن نال وی جڑی اے۔
ساڈا مقصد تہانوں سیٹ نظریے دیاں اوہ بنیادی گلاں سمجھانا اے جیہڑیاں سیٹ
نظریے دے فلسفے دی شروعات لئی چاہیدیاں نیں۔ تے اوس فلسفیانہ لکھت وچ زیادہ تر
ایسے سیٹ نظریاں دی گل اے جیہڑے غیر سیٹ بنیادی شےواں توں \emph{بناں} لکھے
گئے نیں۔ سو ایہ کتاب شاید ہور کم آوے جے اوہ اوس لکھت دے کافی نیڑے رہے۔

\end{document}
